\chapter{Evaluation}
\label{sec:evaluation}

\section{Testumgebung und Datensätze}
Die Evaluation des \textit{multimodalen Sprach-Assistenten} erfolgte auf einem lokalen Entwicklungsrechner unter macOS sowie innerhalb einer isolierten Python-Umgebung (venv).

Zur funktionalen Überprüfung des Systems wurden verschiedene Sprach- und Texteingaben verwendet, die unterschiedliche linguistische Strukturen abbilden. Diese umfassen einfache Sätze, komplexe Mehrsatzstrukturen sowie handlungsorientierte Beschreibungen, um die Robustheit der NLP-Pipeline zu testen.

\section{Performance-Analyse}

\subsection{Speech-to-Text (Whisper)}
Die Spracherkennung mittels OpenAI Whisper zeigte eine hohe Genauigkeit bei klar gesprochenen Eingaben.

\begin{itemize}
    \item Gute Erkennungsrate bei deutschen Standardsätzen
    \item Leichte Genauigkeitsverluste bei stark verrauschten Audioeingaben
    \item Transkriptionszeit abhängig von Modellgröße (base vs. small)
\end{itemize}

Die Verarbeitung erfolgt nahezu in Echtzeit für kurze Audiosequenzen, wodurch eine interaktive Nutzung ermöglicht wird.

\subsection{Natural Language Processing (spaCy)}
Die linguistische Analyse durch spaCy liefert stabile Ergebnisse im Bereich Tokenisierung und Part-of-Speech Tagging.

\begin{itemize}
    \item Verben werden zuverlässig als Handlungen erkannt
    \item Nomen werden konsistent als Objekte extrahiert
    \item Satzstruktur bleibt auch bei längeren Eingaben stabil analysierbar
\end{itemize}

Die Performance der NLP-Komponente ist unabhängig von der Eingabelänge für typische Nutzungsszenarien ausreichend.

\subsection{Bildgenerierung und semantische Abbildung}
Die semantische Umwandlung von Sprache in visuelle Inhalte erfolgt über die Extraktion von Nomen, welche als Suchbegriffe für eine Bilddatenquelle verwendet werden.

\begin{itemize}
    \item Hohe Relevanz bei einfachen Objektbegriffen (z.B. "Hund", "Park")
    \item Eingeschränkte Genauigkeit bei abstrakten Begriffen
    \item Schnelle Bildauslieferung über externe API (Unsplash)
\end{itemize}

\section{Funktionale Ergebnisse}

Die funktionale Evaluation des Gesamtsystems ergab folgende Erkenntnisse:

\begin{itemize}
    \item \textbf{End-to-End Pipeline:} Die vollständige Verarbeitungskette (Audio $\rightarrow$ Text $\rightarrow$ NLP $\rightarrow$ Bild) funktioniert stabil und ohne manuelle Eingriffe.
    
    \item \textbf{Textvisualisierung:} Die Hervorhebung von Verben im Text verbessert die Nachvollziehbarkeit der linguistischen Analyse deutlich.

    \item \textbf{Interaktive Nutzung:} Die Streamlit-Oberfläche ermöglicht eine einfache Bedienung und unmittelbare Rückmeldung der KI-Ergebnisse.

    \item \textbf{Robustheit:} Das System bleibt auch bei fehlerhaften oder unvollständigen Eingaben stabil und liefert interpretierbare Ergebnisse.
\end{itemize}

\section{Zusammenfassende Beurteilung}
Die Evaluation zeigt, dass der entwickelte multimodale Sprach-Assistent die definierten Anforderungen erfüllt. Besonders hervorzuheben ist die erfolgreiche Kombination aus Spracherkennung, linguistischer Analyse und visueller Repräsentation.

Die modulare Architektur erlaubt eine einfache Erweiterung des Systems, beispielsweise durch den Einsatz leistungsfähigerer Sprachmodelle oder generativer Bild-KI. Insgesamt stellt das Projekt eine funktionale KI-Pipeline dar, die Sprache in strukturierte und visuell interpretierbare Informationen überführt.